% --- Archon physics formalization source begin ---
% archon:physics
% archon:covers PhyXMiniProblems/problem_phyx_mini_0996.lean
% archon:source-report reports/phyx_mini/problem_phyx_mini_0996.source.json

\chapter{Physics problem phyx\_mini\_0996}
\label{ch:PhyXMiniProblems_problem_phyx_mini_0996}

\paragraph{Problem source.}
\#\# Physical scenario

Figure shows an electric dipole in a uniform electric field of magnitude \textbackslash{}( 5.0 \textbackslash{}times 10\textasciicircum{}5\textbackslash{} \textbackslash{}text\{N/C\} \textbackslash{}) that is directed parallel to the plane of the figure. The charges are \textbackslash{}( \textbackslash{}pm 1.6 \textbackslash{}times 10\textasciicircum{}\{-19\}\textbackslash{} \textbackslash{}text\{C\} \textbackslash{}); both lie in the plane and are separated by \textbackslash{}( 0.125\textbackslash{} \textbackslash{}text\{m\} = 0.125 \textbackslash{}times 10\textasciicircum{}\{-9\}\textbackslash{} \textbackslash{}text\{m\} \textbackslash{}).

\#\# Question

Find the magnitude of the torque.

\#\# Answer choices

- A: A: 2.0 \textbackslash{}times 10\textasciicircum{}\{-29\}\textbackslash{}text\{N\} \textbackslash{}cdot \textbackslash{}text\{m\}

- B: B: 5.7 \textbackslash{}times 10\textasciicircum{}\{-24\} \textbackslash{}

- C: C: 8.2 \textbackslash{}times 10\textasciicircum{}\{-24\}\textbackslash{}text\{J\}

- D: D: 5.7 \textbackslash{}times 10\textasciicircum{}\{-24\}\textbackslash{}text\{N/C\}

\#\# Figure caption (auxiliary; use the image as primary evidence)

The image contains a diagram with several vector arrows emanating from a central point, which is marked with a red dot. 

1. **Vectors**:
   - A teal vector labeled \textbackslash{}(\textbackslash{}overrightarrow\{p\}\textbackslash{}) points downward and to the left.
   - A magenta vector labeled \textbackslash{}(\textbackslash{}overrightarrow\{E\}\textbackslash{}) points horizontally to the right.
   - A smaller arrow labeled \textbackslash{}(\textbackslash{}overrightarrow\{\textbackslash{}tau\}\textbackslash{}) is positioned above the central point, indicating the direction between the vectors.

2. **Angle and Arc**:
   - There is a black curved arc marking the angle between \textbackslash{}(\textbackslash{}overrightarrow\{p\}\textbackslash{}) and \textbackslash{}(\textbackslash{}overrightarrow\{E\}\textbackslash{}).
   - The angle is labeled as \textbackslash{}(145\textasciicircum{}\textbackslash{}circ\textbackslash{}) near the arc, highlighting the degree measure between the two main vectors.

The vectors illustrate relationships such as direction and magnitude, and the angle suggests an orientation between the vectors \textbackslash{}(\textbackslash{}overrightarrow\{p\}\textbackslash{}) and \textbackslash{}(\textbackslash{}overrightarrow\{E\}\textbackslash{}).

\#\# Dataset metadata

Electrostatics; Physical Model Grounding Reasoning, Spatial Relation Reasoning

\paragraph{Recorded answer/context.}
B: B: 5.7 \textbackslash{}times 10\textasciicircum{}\{-24\} \textbackslash{}

\paragraph{Figure/image path.}
/root/proposal\_for\_physic/science-mango/phyx\_mini\_run/phyx\_data/test\_image/996.png

\paragraph{Formalization target.}
create a compiling Lean file with sorry bodies at `PhyXMiniProblems/problem\_phyx\_mini\_0996.lean`. The Lean declarations must preserve the physical quantities, dimensions or dimensional roles, figure labels, governing-law hypotheses, and final relation expressed by this problem.
Use Mathlib/Physlib names found through LeanExplore where available. If a physics API is missing, introduce faithful local abstractions rather than scalar placeholder aliases.

\begin{theorem}[Physics formalization target]
\label{thm:physics:phyx_mini_0996:target}
\lean{PhyXMiniProblems.ProblemPhyXMini0996.problem_phyx_mini_0996}
\uses{def:physics:phyx-mini-0996:phyxminiproblems-problemphyxmini0996-torquemagnitudeinnewtonmeters, def:physics:phyx-mini-0996:phyxminiproblems-problemphyxmini0996-electricdipoletorquesetup, def:physics:phyx-mini-0996:phyxminiproblems-problemphyxmini0996-matcheswrittenscenario, def:physics:phyx-mini-0996:phyxminiproblems-problemphyxmini0996-matchesproblemreadouts, def:physics:phyx-mini-0996:phyxminiproblems-problemphyxmini0996-matchesprimaryfigure, def:physics:phyx-mini-0996:phyxminiproblems-problemphyxmini0996-haselectricdipolegeometry, def:physics:phyx-mini-0996:phyxminiproblems-problemphyxmini0996-representsuniformelectricfield, def:physics:phyx-mini-0996:phyxminiproblems-problemphyxmini0996-obeyselectricdipoletorquelaws, def:physics:phyx-mini-0996:phyxminiproblems-problemphyxmini0996-roundstoatresolution, def:physics:phyx-mini-0996:phyxminiproblems-problemphyxmini0996-isuniquematchingtorquechoice}
The assigned autoformalize agent should translate this physics problem into Lean declarations in the covered file, with theorem and lemma proof bodies written as `by sorry`.
\end{theorem}
\begin{proof}
This is an autoformalization task, not a proof task. Produce statements that can later be proved by the physics prover without weakening the physical model.
\end{proof}

% --- Archon declaration topology begin ---
\section{Declaration topology}
\begin{definition}[electric Field Strength Dimension]
  \label{def:physics:phyx-mini-0996:phyxminiproblems-problemphyxmini0996-electricfieldstrengthdimension}
  \lean{PhyXMiniProblems.ProblemPhyXMini0996.electricFieldStrengthDimension}
  The dimension M L T⁻² C⁻¹ of electric-field strength
\end{definition}
\begin{proof}
  This is the declared model definition of electric Field Strength Dimension.
\end{proof}
\begin{definition}[electric Dipole Moment Dimension]
  \label{def:physics:phyx-mini-0996:phyxminiproblems-problemphyxmini0996-electricdipolemomentdimension}
  \lean{PhyXMiniProblems.ProblemPhyXMini0996.electricDipoleMomentDimension}
  The dimension C L of electric dipole moment
\end{definition}
\begin{proof}
  This is the declared model definition of electric Dipole Moment Dimension.
\end{proof}
\begin{definition}[torque Dimension]
  \label{def:physics:phyx-mini-0996:phyxminiproblems-problemphyxmini0996-torquedimension}
  \lean{PhyXMiniProblems.ProblemPhyXMini0996.torqueDimension}
  The dimension M L² T⁻² of torque
\end{definition}
\begin{proof}
  This is the declared model definition of torque Dimension.
\end{proof}
\begin{definition}[Spatial Vector]
  \label{def:physics:phyx-mini-0996:phyxminiproblems-problemphyxmini0996-spatialvector}
  \lean{PhyXMiniProblems.ProblemPhyXMini0996.SpatialVector}
  Three-dimensional vectors used at coherent-SI readout boundaries
\end{definition}
\begin{proof}
  This is the declared model definition of Spatial Vector.
\end{proof}
\begin{definition}[Signed Charge Quantity]
  \label{def:physics:phyx-mini-0996:phyxminiproblems-problemphyxmini0996-signedchargequantity}
  \lean{PhyXMiniProblems.ProblemPhyXMini0996.SignedChargeQuantity}
  A signed, unit-independent electric charge
\end{definition}
\begin{proof}
  This is the declared model definition of Signed Charge Quantity.
\end{proof}
\begin{definition}[Length Quantity]
  \label{def:physics:phyx-mini-0996:phyxminiproblems-problemphyxmini0996-lengthquantity}
  \lean{PhyXMiniProblems.ProblemPhyXMini0996.LengthQuantity}
  A nonnegative, unit-independent physical length
\end{definition}
\begin{proof}
  This is the declared model definition of Length Quantity.
\end{proof}
\begin{definition}[Position Quantity]
  \label{def:physics:phyx-mini-0996:phyxminiproblems-problemphyxmini0996-positionquantity}
  \lean{PhyXMiniProblems.ProblemPhyXMini0996.PositionQuantity}
  \uses{def:physics:phyx-mini-0996:phyxminiproblems-problemphyxmini0996-spatialvector}
  A unit-independent physical position vector
\end{definition}
\begin{proof}
  This is the declared model definition of Position Quantity.
\end{proof}
\begin{definition}[Electric Field Strength Quantity]
  \label{def:physics:phyx-mini-0996:phyxminiproblems-problemphyxmini0996-electricfieldstrengthquantity}
  \lean{PhyXMiniProblems.ProblemPhyXMini0996.ElectricFieldStrengthQuantity}
  \uses{def:physics:phyx-mini-0996:phyxminiproblems-problemphyxmini0996-electricfieldstrengthdimension}
  A nonnegative, unit-independent electric-field magnitude
\end{definition}
\begin{proof}
  This is the declared model definition of Electric Field Strength Quantity.
\end{proof}
\begin{definition}[Electric Field Vector Quantity]
  \label{def:physics:phyx-mini-0996:phyxminiproblems-problemphyxmini0996-electricfieldvectorquantity}
  \lean{PhyXMiniProblems.ProblemPhyXMini0996.ElectricFieldVectorQuantity}
  \uses{def:physics:phyx-mini-0996:phyxminiproblems-problemphyxmini0996-electricfieldstrengthdimension, def:physics:phyx-mini-0996:phyxminiproblems-problemphyxmini0996-spatialvector}
  A unit-independent electric-field vector
\end{definition}
\begin{proof}
  This is the declared model definition of Electric Field Vector Quantity.
\end{proof}
\begin{definition}[Electric Dipole Moment Quantity]
  \label{def:physics:phyx-mini-0996:phyxminiproblems-problemphyxmini0996-electricdipolemomentquantity}
  \lean{PhyXMiniProblems.ProblemPhyXMini0996.ElectricDipoleMomentQuantity}
  \uses{def:physics:phyx-mini-0996:phyxminiproblems-problemphyxmini0996-electricdipolemomentdimension, def:physics:phyx-mini-0996:phyxminiproblems-problemphyxmini0996-spatialvector}
  A unit-independent electric-dipole-moment vector
\end{definition}
\begin{proof}
  This is the declared model definition of Electric Dipole Moment Quantity.
\end{proof}
\begin{definition}[Torque Vector Quantity]
  \label{def:physics:phyx-mini-0996:phyxminiproblems-problemphyxmini0996-torquevectorquantity}
  \lean{PhyXMiniProblems.ProblemPhyXMini0996.TorqueVectorQuantity}
  \uses{def:physics:phyx-mini-0996:phyxminiproblems-problemphyxmini0996-torquedimension, def:physics:phyx-mini-0996:phyxminiproblems-problemphyxmini0996-spatialvector}
  A unit-independent electric-torque vector
\end{definition}
\begin{proof}
  This is the declared model definition of Torque Vector Quantity.
\end{proof}
\begin{definition}[Torque Magnitude Quantity]
  \label{def:physics:phyx-mini-0996:phyxminiproblems-problemphyxmini0996-torquemagnitudequantity}
  \lean{PhyXMiniProblems.ProblemPhyXMini0996.TorqueMagnitudeQuantity}
  \uses{def:physics:phyx-mini-0996:phyxminiproblems-problemphyxmini0996-torquedimension}
  A nonnegative, unit-independent torque magnitude
\end{definition}
\begin{proof}
  This is the declared model definition of Torque Magnitude Quantity.
\end{proof}
\begin{definition}[signed Charge In Coulombs]
  \label{def:physics:phyx-mini-0996:phyxminiproblems-problemphyxmini0996-signedchargeincoulombs}
  \lean{PhyXMiniProblems.ProblemPhyXMini0996.signedChargeInCoulombs}
  \uses{def:physics:phyx-mini-0996:phyxminiproblems-problemphyxmini0996-signedchargequantity}
  Read a signed charge in coherent-SI coulombs
\end{definition}
\begin{proof}
  This is the declared model definition of signed Charge In Coulombs.
\end{proof}
\begin{definition}[length In Meters]
  \label{def:physics:phyx-mini-0996:phyxminiproblems-problemphyxmini0996-lengthinmeters}
  \lean{PhyXMiniProblems.ProblemPhyXMini0996.lengthInMeters}
  \uses{def:physics:phyx-mini-0996:phyxminiproblems-problemphyxmini0996-lengthquantity}
  Read a length in coherent-SI metres
\end{definition}
\begin{proof}
  This is the declared model definition of length In Meters.
\end{proof}
\begin{definition}[position Vector In Meters]
  \label{def:physics:phyx-mini-0996:phyxminiproblems-problemphyxmini0996-positionvectorinmeters}
  \lean{PhyXMiniProblems.ProblemPhyXMini0996.positionVectorInMeters}
  \uses{def:physics:phyx-mini-0996:phyxminiproblems-problemphyxmini0996-spatialvector, def:physics:phyx-mini-0996:phyxminiproblems-problemphyxmini0996-positionquantity}
  Read a physical position vector in coherent-SI metres
\end{definition}
\begin{proof}
  This is the declared model definition of position Vector In Meters.
\end{proof}
\begin{definition}[electric Field Strength In Newtons Per Coulomb]
  \label{def:physics:phyx-mini-0996:phyxminiproblems-problemphyxmini0996-electricfieldstrengthinnewtonspercoulomb}
  \lean{PhyXMiniProblems.ProblemPhyXMini0996.electricFieldStrengthInNewtonsPerCoulomb}
  \uses{def:physics:phyx-mini-0996:phyxminiproblems-problemphyxmini0996-electricfieldstrengthquantity}
  Read an electric-field magnitude in coherent-SI newtons per coulomb
\end{definition}
\begin{proof}
  This is the declared model definition of electric Field Strength In Newtons Per Coulomb.
\end{proof}
\begin{definition}[electric Field Vector In Newtons Per Coulomb]
  \label{def:physics:phyx-mini-0996:phyxminiproblems-problemphyxmini0996-electricfieldvectorinnewtonspercoulomb}
  \lean{PhyXMiniProblems.ProblemPhyXMini0996.electricFieldVectorInNewtonsPerCoulomb}
  \uses{def:physics:phyx-mini-0996:phyxminiproblems-problemphyxmini0996-spatialvector, def:physics:phyx-mini-0996:phyxminiproblems-problemphyxmini0996-electricfieldvectorquantity}
  Read an electric-field vector in coherent-SI newtons per coulomb
\end{definition}
\begin{proof}
  This is the declared model definition of electric Field Vector In Newtons Per Coulomb.
\end{proof}
\begin{definition}[dipole Moment Vector In Coulomb Meters]
  \label{def:physics:phyx-mini-0996:phyxminiproblems-problemphyxmini0996-dipolemomentvectorincoulombmeters}
  \lean{PhyXMiniProblems.ProblemPhyXMini0996.dipoleMomentVectorInCoulombMeters}
  \uses{def:physics:phyx-mini-0996:phyxminiproblems-problemphyxmini0996-spatialvector, def:physics:phyx-mini-0996:phyxminiproblems-problemphyxmini0996-electricdipolemomentquantity}
  Read a dipole-moment vector in coherent-SI coulomb metres
\end{definition}
\begin{proof}
  This is the declared model definition of dipole Moment Vector In Coulomb Meters.
\end{proof}
\begin{definition}[torque Vector In Newton Meters]
  \label{def:physics:phyx-mini-0996:phyxminiproblems-problemphyxmini0996-torquevectorinnewtonmeters}
  \lean{PhyXMiniProblems.ProblemPhyXMini0996.torqueVectorInNewtonMeters}
  \uses{def:physics:phyx-mini-0996:phyxminiproblems-problemphyxmini0996-spatialvector, def:physics:phyx-mini-0996:phyxminiproblems-problemphyxmini0996-torquevectorquantity}
  Read an electric-torque vector in coherent-SI newton metres
\end{definition}
\begin{proof}
  This is the declared model definition of torque Vector In Newton Meters.
\end{proof}
\begin{definition}[torque Magnitude In Newton Meters]
  \label{def:physics:phyx-mini-0996:phyxminiproblems-problemphyxmini0996-torquemagnitudeinnewtonmeters}
  \lean{PhyXMiniProblems.ProblemPhyXMini0996.torqueMagnitudeInNewtonMeters}
  \uses{def:physics:phyx-mini-0996:phyxminiproblems-problemphyxmini0996-torquemagnitudequantity}
  Read a torque magnitude in coherent-SI newton metres
\end{definition}
\begin{proof}
  This is the declared model definition of torque Magnitude In Newton Meters.
\end{proof}
\begin{definition}[degrees To Radians]
  \label{def:physics:phyx-mini-0996:phyxminiproblems-problemphyxmini0996-degreestoradians}
  \lean{PhyXMiniProblems.ProblemPhyXMini0996.degreesToRadians}
  Convert an angle measured in degrees to a dimensionless radian value
\end{definition}
\begin{proof}
  This is the declared model definition of degrees To Radians.
\end{proof}
\begin{definition}[Charge Label]
  \label{def:physics:phyx-mini-0996:phyxminiproblems-problemphyxmini0996-chargelabel}
  \lean{PhyXMiniProblems.ProblemPhyXMini0996.ChargeLabel}
  The signs of the two point charges making up the dipole
\end{definition}
\begin{proof}
  This is the declared model definition of Charge Label.
\end{proof}
\begin{definition}[Figure Vector Label]
  \label{def:physics:phyx-mini-0996:phyxminiproblems-problemphyxmini0996-figurevectorlabel}
  \lean{PhyXMiniProblems.ProblemPhyXMini0996.FigureVectorLabel}
  The three vector symbols printed in the supplied image
\end{definition}
\begin{proof}
  This is the declared model definition of Figure Vector Label.
\end{proof}
\begin{definition}[Figure Arrow Direction]
  \label{def:physics:phyx-mini-0996:phyxminiproblems-problemphyxmini0996-figurearrowdirection}
  \lean{PhyXMiniProblems.ProblemPhyXMini0996.FigureArrowDirection}
  Qualitative arrow directions visible in the plane of the image
\end{definition}
\begin{proof}
  This is the declared model definition of Figure Arrow Direction.
\end{proof}
\begin{definition}[Dipole Model]
  \label{def:physics:phyx-mini-0996:phyxminiproblems-problemphyxmini0996-dipolemodel}
  \lean{PhyXMiniProblems.ProblemPhyXMini0996.DipoleModel}
  Idealization of the two charges as an electric dipole
\end{definition}
\begin{proof}
  This is the declared model definition of Dipole Model.
\end{proof}
\begin{definition}[Electric Field Spatial Model]
  \label{def:physics:phyx-mini-0996:phyxminiproblems-problemphyxmini0996-electricfieldspatialmodel}
  \lean{PhyXMiniProblems.ProblemPhyXMini0996.ElectricFieldSpatialModel}
  Spatial dependence assigned to the external electric field
\end{definition}
\begin{proof}
  This is the declared model definition of Electric Field Spatial Model.
\end{proof}
\begin{definition}[Figure Plane Relation]
  \label{def:physics:phyx-mini-0996:phyxminiproblems-problemphyxmini0996-figureplanerelation}
  \lean{PhyXMiniProblems.ProblemPhyXMini0996.FigurePlaneRelation}
  Relation of a vector to the plane containing the primary diagram
\end{definition}
\begin{proof}
  This is the declared model definition of Figure Plane Relation.
\end{proof}
\begin{definition}[Electric Dipole Torque Figure]
  \label{def:physics:phyx-mini-0996:phyxminiproblems-problemphyxmini0996-electricdipoletorquefigure}
  \lean{PhyXMiniProblems.ProblemPhyXMini0996.ElectricDipoleTorqueFigure}
  \uses{def:physics:phyx-mini-0996:phyxminiproblems-problemphyxmini0996-figurevectorlabel, def:physics:phyx-mini-0996:phyxminiproblems-problemphyxmini0996-figurearrowdirection}
  Literal graphical evidence from 996.png. It records symbols, directions, the red common origin, and the angle arc, but no torque-magnitude answer
\end{definition}
\begin{proof}
  This is the declared model definition of Electric Dipole Torque Figure.
\end{proof}
\begin{definition}[Electric Dipole Torque Setup]
  \label{def:physics:phyx-mini-0996:phyxminiproblems-problemphyxmini0996-electricdipoletorquesetup}
  \lean{PhyXMiniProblems.ProblemPhyXMini0996.ElectricDipoleTorqueSetup}
  \uses{def:physics:phyx-mini-0996:phyxminiproblems-problemphyxmini0996-spatialvector, def:physics:phyx-mini-0996:phyxminiproblems-problemphyxmini0996-signedchargequantity, def:physics:phyx-mini-0996:phyxminiproblems-problemphyxmini0996-lengthquantity, def:physics:phyx-mini-0996:phyxminiproblems-problemphyxmini0996-positionquantity, def:physics:phyx-mini-0996:phyxminiproblems-problemphyxmini0996-electricfieldstrengthquantity, def:physics:phyx-mini-0996:phyxminiproblems-problemphyxmini0996-electricfieldvectorquantity, def:physics:phyx-mini-0996:phyxminiproblems-problemphyxmini0996-electricdipolemomentquantity, def:physics:phyx-mini-0996:phyxminiproblems-problemphyxmini0996-torquevectorquantity, def:physics:phyx-mini-0996:phyxminiproblems-problemphyxmini0996-torquemagnitudequantity, def:physics:phyx-mini-0996:phyxminiproblems-problemphyxmini0996-chargelabel, def:physics:phyx-mini-0996:phyxminiproblems-problemphyxmini0996-dipolemodel, def:physics:phyx-mini-0996:phyxminiproblems-problemphyxmini0996-electricfieldspatialmodel, def:physics:phyx-mini-0996:phyxminiproblems-problemphyxmini0996-figureplanerelation, def:physics:phyx-mini-0996:phyxminiproblems-problemphyxmini0996-electricdipoletorquefigure}
  The charges, positions, field, dipole moment, torque vector, and torque magnitude are independent physical data. Their relationships occur only in the scenario, geometry, and governing-law predicates below
\end{definition}
\begin{proof}
  This is the declared model definition of Electric Dipole Torque Setup.
\end{proof}
\begin{definition}[charge Separation Vector In Meters]
  \label{def:physics:phyx-mini-0996:phyxminiproblems-problemphyxmini0996-chargeseparationvectorinmeters}
  \lean{PhyXMiniProblems.ProblemPhyXMini0996.chargeSeparationVectorInMeters}
  \uses{def:physics:phyx-mini-0996:phyxminiproblems-problemphyxmini0996-spatialvector, def:physics:phyx-mini-0996:phyxminiproblems-problemphyxmini0996-positionvectorinmeters, def:physics:phyx-mini-0996:phyxminiproblems-problemphyxmini0996-electricdipoletorquesetup}
  Displacement from the negative charge to the positive charge, in metres
\end{definition}
\begin{proof}
  This is the declared model definition of charge Separation Vector In Meters.
\end{proof}
\begin{definition}[Points Rightward In Figure]
  \label{def:physics:phyx-mini-0996:phyxminiproblems-problemphyxmini0996-pointsrightwardinfigure}
  \lean{PhyXMiniProblems.ProblemPhyXMini0996.PointsRightwardInFigure}
  \uses{def:physics:phyx-mini-0996:phyxminiproblems-problemphyxmini0996-spatialvector}
  A vector points strictly right in the chosen in-page coordinates
\end{definition}
\begin{proof}
  This is the declared model definition of Points Rightward In Figure.
\end{proof}
\begin{definition}[Points Downward Left In Figure]
  \label{def:physics:phyx-mini-0996:phyxminiproblems-problemphyxmini0996-pointsdownwardleftinfigure}
  \lean{PhyXMiniProblems.ProblemPhyXMini0996.PointsDownwardLeftInFigure}
  \uses{def:physics:phyx-mini-0996:phyxminiproblems-problemphyxmini0996-spatialvector}
  A vector points down and left in the chosen in-page coordinates
\end{definition}
\begin{proof}
  This is the declared model definition of Points Downward Left In Figure.
\end{proof}
\begin{definition}[Matches Written Scenario]
  \label{def:physics:phyx-mini-0996:phyxminiproblems-problemphyxmini0996-matcheswrittenscenario}
  \lean{PhyXMiniProblems.ProblemPhyXMini0996.MatchesWrittenScenario}
  \uses{def:physics:phyx-mini-0996:phyxminiproblems-problemphyxmini0996-electricdipoletorquesetup}
  Qualitative facts stated in the written physical scenario
\end{definition}
\begin{proof}
  This is the declared model definition of Matches Written Scenario.
\end{proof}
\begin{definition}[Matches Problem Readouts]
  \label{def:physics:phyx-mini-0996:phyxminiproblems-problemphyxmini0996-matchesproblemreadouts}
  \lean{PhyXMiniProblems.ProblemPhyXMini0996.MatchesProblemReadouts}
  \uses{def:physics:phyx-mini-0996:phyxminiproblems-problemphyxmini0996-signedchargeincoulombs, def:physics:phyx-mini-0996:phyxminiproblems-problemphyxmini0996-lengthinmeters, def:physics:phyx-mini-0996:phyxminiproblems-problemphyxmini0996-electricfieldstrengthinnewtonspercoulomb, def:physics:phyx-mini-0996:phyxminiproblems-problemphyxmini0996-electricdipoletorquesetup}
  Numerical problem data. The separation is read as 0.125 nm = 0.125 × 10⁻⁹ m. No torque value occurs in these premises
\end{definition}
\begin{proof}
  This is the declared model definition of Matches Problem Readouts.
\end{proof}
\begin{definition}[Matches Primary Figure]
  \label{def:physics:phyx-mini-0996:phyxminiproblems-problemphyxmini0996-matchesprimaryfigure}
  \lean{PhyXMiniProblems.ProblemPhyXMini0996.MatchesPrimaryFigure}
  \uses{def:physics:phyx-mini-0996:phyxminiproblems-problemphyxmini0996-degreestoradians, def:physics:phyx-mini-0996:phyxminiproblems-problemphyxmini0996-electricdipoletorquesetup}
  Primary-image evidence and its calibration against the stored physical angle. The raster does not supply the requested torque magnitude
\end{definition}
\begin{proof}
  This is the declared model definition of Matches Primary Figure.
\end{proof}
\begin{definition}[Has Electric Dipole Geometry]
  \label{def:physics:phyx-mini-0996:phyxminiproblems-problemphyxmini0996-haselectricdipolegeometry}
  \lean{PhyXMiniProblems.ProblemPhyXMini0996.HasElectricDipoleGeometry}
  \uses{def:physics:phyx-mini-0996:phyxminiproblems-problemphyxmini0996-lengthinmeters, def:physics:phyx-mini-0996:phyxminiproblems-problemphyxmini0996-positionvectorinmeters, def:physics:phyx-mini-0996:phyxminiproblems-problemphyxmini0996-electricfieldvectorinnewtonspercoulomb, def:physics:phyx-mini-0996:phyxminiproblems-problemphyxmini0996-dipolemomentvectorincoulombmeters, def:physics:phyx-mini-0996:phyxminiproblems-problemphyxmini0996-electricdipoletorquesetup, def:physics:phyx-mini-0996:phyxminiproblems-problemphyxmini0996-chargeseparationvectorinmeters, def:physics:phyx-mini-0996:phyxminiproblems-problemphyxmini0996-pointsrightwardinfigure, def:physics:phyx-mini-0996:phyxminiproblems-problemphyxmini0996-pointsdownwardleftinfigure}
  Geometric meaning of the charge positions, in-page vectors, and marked angle
\end{definition}
\begin{proof}
  This is the declared model definition of Has Electric Dipole Geometry.
\end{proof}
\begin{definition}[Represents Uniform Electric Field]
  \label{def:physics:phyx-mini-0996:phyxminiproblems-problemphyxmini0996-representsuniformelectricfield}
  \lean{PhyXMiniProblems.ProblemPhyXMini0996.RepresentsUniformElectricField}
  \uses{def:physics:phyx-mini-0996:phyxminiproblems-problemphyxmini0996-electricfieldstrengthinnewtonspercoulomb, def:physics:phyx-mini-0996:phyxminiproblems-problemphyxmini0996-electricfieldvectorinnewtonspercoulomb, def:physics:phyx-mini-0996:phyxminiproblems-problemphyxmini0996-electricdipoletorquesetup}
  Physlib's spacetime-dependent field is specialized to the stated static, uniform vector. Its norm agrees with the dimensionful scalar field strength
\end{definition}
\begin{proof}
  This is the declared model definition of Represents Uniform Electric Field.
\end{proof}
\begin{definition}[Obeys Electric Dipole Torque Laws]
  \label{def:physics:phyx-mini-0996:phyxminiproblems-problemphyxmini0996-obeyselectricdipoletorquelaws}
  \lean{PhyXMiniProblems.ProblemPhyXMini0996.ObeysElectricDipoleTorqueLaws}
  \uses{def:physics:phyx-mini-0996:phyxminiproblems-problemphyxmini0996-signedchargeincoulombs, def:physics:phyx-mini-0996:phyxminiproblems-problemphyxmini0996-electricfieldvectorinnewtonspercoulomb, def:physics:phyx-mini-0996:phyxminiproblems-problemphyxmini0996-dipolemomentvectorincoulombmeters, def:physics:phyx-mini-0996:phyxminiproblems-problemphyxmini0996-torquevectorinnewtonmeters, def:physics:phyx-mini-0996:phyxminiproblems-problemphyxmini0996-torquemagnitudeinnewtonmeters, def:physics:phyx-mini-0996:phyxminiproblems-problemphyxmini0996-electricdipoletorquesetup, def:physics:phyx-mini-0996:phyxminiproblems-problemphyxmini0996-chargeseparationvectorinmeters}
  General point-dipole and electric-torque laws. They state q₋ = -q₊, p = q₊ (r₊ - r₋), τ = p × E, and that the measured torque magnitude is the norm of the torque vector. They contain no requested numerical answer
\end{definition}
\begin{proof}
  This is the declared model definition of Obeys Electric Dipole Torque Laws.
\end{proof}
\begin{definition}[Answer Choice]
  \label{def:physics:phyx-mini-0996:phyxminiproblems-problemphyxmini0996-answerchoice}
  \lean{PhyXMiniProblems.ProblemPhyXMini0996.AnswerChoice}
  Labels of the four multiple-choice answers in the source
\end{definition}
\begin{proof}
  This is the declared model definition of Answer Choice.
\end{proof}
\begin{definition}[Displayed Unit]
  \label{def:physics:phyx-mini-0996:phyxminiproblems-problemphyxmini0996-displayedunit}
  \lean{PhyXMiniProblems.ProblemPhyXMini0996.DisplayedUnit}
  Unit text printed beside each answer choice
\end{definition}
\begin{proof}
  This is the declared model definition of Displayed Unit.
\end{proof}
\begin{definition}[displayed Magnitude SI]
  \label{def:physics:phyx-mini-0996:phyxminiproblems-problemphyxmini0996-answerchoice-displayedmagnitudesi}
  \lean{PhyXMiniProblems.ProblemPhyXMini0996.AnswerChoice.displayedMagnitudeSI}
  \uses{def:physics:phyx-mini-0996:phyxminiproblems-problemphyxmini0996-answerchoice}
  Numerical value printed by each answer, interpreted as an SI scalar
\end{definition}
\begin{proof}
  This is the declared model definition of displayed Magnitude SI.
\end{proof}
\begin{definition}[displayed Unit]
  \label{def:physics:phyx-mini-0996:phyxminiproblems-problemphyxmini0996-answerchoice-displayedunit}
  \lean{PhyXMiniProblems.ProblemPhyXMini0996.AnswerChoice.displayedUnit}
  \uses{def:physics:phyx-mini-0996:phyxminiproblems-problemphyxmini0996-answerchoice, def:physics:phyx-mini-0996:phyxminiproblems-problemphyxmini0996-displayedunit}
  Unit text attached to each displayed choice in the source
\end{definition}
\begin{proof}
  This is the declared model definition of displayed Unit.
\end{proof}
\begin{definition}[interpreted Unit For Torque Question]
  \label{def:physics:phyx-mini-0996:phyxminiproblems-problemphyxmini0996-answerchoice-interpretedunitfortorquequestion}
  \lean{PhyXMiniProblems.ProblemPhyXMini0996.AnswerChoice.interpretedUnitForTorqueQuestion}
  \uses{def:physics:phyx-mini-0996:phyxminiproblems-problemphyxmini0996-answerchoice, def:physics:phyx-mini-0996:phyxminiproblems-problemphyxmini0996-displayedunit}
  The source omits B's unit. Since the question explicitly requests a torque, the physically interpreted unit for that option is newton metres. Keeping this separate from displayedUnit preserves both the literal source and the unit needed to understand the proposed answer
\end{definition}
\begin{proof}
  This is the declared model definition of interpreted Unit For Torque Question.
\end{proof}
\begin{definition}[display Resolution]
  \label{def:physics:phyx-mini-0996:phyxminiproblems-problemphyxmini0996-answerchoice-displayresolution}
  \lean{PhyXMiniProblems.ProblemPhyXMini0996.AnswerChoice.displayResolution}
  \uses{def:physics:phyx-mini-0996:phyxminiproblems-problemphyxmini0996-answerchoice}
  One unit in the last displayed decimal place of each answer
\end{definition}
\begin{proof}
  This is the declared model definition of display Resolution.
\end{proof}
\begin{definition}[is Torque Compatible]
  \label{def:physics:phyx-mini-0996:phyxminiproblems-problemphyxmini0996-displayedunit-istorquecompatible}
  \lean{PhyXMiniProblems.ProblemPhyXMini0996.DisplayedUnit.isTorqueCompatible}
  \uses{def:physics:phyx-mini-0996:phyxminiproblems-problemphyxmini0996-displayedunit}
  Whether an explicit unit is dimensionally compatible with torque
\end{definition}
\begin{proof}
  This is the declared model definition of is Torque Compatible.
\end{proof}
\begin{definition}[Rounds To At Resolution]
  \label{def:physics:phyx-mini-0996:phyxminiproblems-problemphyxmini0996-roundstoatresolution}
  \lean{PhyXMiniProblems.ProblemPhyXMini0996.RoundsToAtResolution}
  An actual scalar rounds to a displayed scalar at the stated resolution
\end{definition}
\begin{proof}
  This is the declared model definition of Rounds To At Resolution.
\end{proof}
\begin{definition}[Matches Displayed Torque Choice]
  \label{def:physics:phyx-mini-0996:phyxminiproblems-problemphyxmini0996-matchesdisplayedtorquechoice}
  \lean{PhyXMiniProblems.ProblemPhyXMini0996.MatchesDisplayedTorqueChoice}
  \uses{def:physics:phyx-mini-0996:phyxminiproblems-problemphyxmini0996-torquemagnitudeinnewtonmeters, def:physics:phyx-mini-0996:phyxminiproblems-problemphyxmini0996-electricdipoletorquesetup, def:physics:phyx-mini-0996:phyxminiproblems-problemphyxmini0996-answerchoice, def:physics:phyx-mini-0996:phyxminiproblems-problemphyxmini0996-roundstoatresolution}
  A choice has torque-compatible units and its displayed number matches
\end{definition}
\begin{proof}
  This is the declared model definition of Matches Displayed Torque Choice.
\end{proof}
\begin{definition}[Is Unique Matching Torque Choice]
  \label{def:physics:phyx-mini-0996:phyxminiproblems-problemphyxmini0996-isuniquematchingtorquechoice}
  \lean{PhyXMiniProblems.ProblemPhyXMini0996.IsUniqueMatchingTorqueChoice}
  \uses{def:physics:phyx-mini-0996:phyxminiproblems-problemphyxmini0996-electricdipoletorquesetup, def:physics:phyx-mini-0996:phyxminiproblems-problemphyxmini0996-answerchoice, def:physics:phyx-mini-0996:phyxminiproblems-problemphyxmini0996-matchesdisplayedtorquechoice}
  A displayed choice is the unique match for the modeled torque magnitude
\end{definition}
\begin{proof}
  This is the declared model definition of Is Unique Matching Torque Choice.
\end{proof}
\begin{definition}[recorded Dataset Answer]
  \label{def:physics:phyx-mini-0996:phyxminiproblems-problemphyxmini0996-recordeddatasetanswer}
  \lean{PhyXMiniProblems.ProblemPhyXMini0996.recordedDatasetAnswer}
  \uses{def:physics:phyx-mini-0996:phyxminiproblems-problemphyxmini0996-answerchoice}
  Dataset metadata recording the supplied answer key; not a theorem premise
\end{definition}
\begin{proof}
  This is the declared model definition of recorded Dataset Answer.
\end{proof}
\begin{lemma}[electric Dipole Torque Magnitude formula]
  \label{lem:physics:phyx-mini-0996:phyxminiproblems-problemphyxmini0996-electricdipoletorquemagnitude-formula}
  \lean{PhyXMiniProblems.ProblemPhyXMini0996.electricDipoleTorqueMagnitude_formula}
  \uses{def:physics:phyx-mini-0996:phyxminiproblems-problemphyxmini0996-signedchargeincoulombs, def:physics:phyx-mini-0996:phyxminiproblems-problemphyxmini0996-lengthinmeters, def:physics:phyx-mini-0996:phyxminiproblems-problemphyxmini0996-electricfieldstrengthinnewtonspercoulomb, def:physics:phyx-mini-0996:phyxminiproblems-problemphyxmini0996-torquemagnitudeinnewtonmeters, def:physics:phyx-mini-0996:phyxminiproblems-problemphyxmini0996-electricdipoletorquesetup, def:physics:phyx-mini-0996:phyxminiproblems-problemphyxmini0996-haselectricdipolegeometry, def:physics:phyx-mini-0996:phyxminiproblems-problemphyxmini0996-representsuniformelectricfield, def:physics:phyx-mini-0996:phyxminiproblems-problemphyxmini0996-obeyselectricdipoletorquelaws}
  The exact general magnitude relation obtained from τ = p × E and the angle between p and E. This relation is derived, not stored in a law field
\end{lemma}
\begin{proof}
  The conclusion follows from the governing relations and figure readouts named in the statement of electric Dipole Torque Magnitude formula.
\end{proof}
% --- Archon declaration topology end ---
% --- Archon physics formalization source end ---
